\section{Introduction}

The integration of the Internet of Medical Things (IoMT) in hospitals, particularly in emergency departments, offers new opportunities for accelerating automated triage and improving early disease detection \cite{esteva2019}. With this kind of environment, diagnostic decisions often rely on the combination of heterogeneous information. Those include vital signs, clinical observations, patient history, laboratory biomarkers, and specialty-specific examinations. To train robust machine learning models for this type of decision support, large amounts of multimodal medical data recreating real-world settings are required.

However, access to real-world patient data remains highly restricted because of privacy and legal constraints. Regulations such as GDPR and HIPAA impose strict limitations on the collection, storage, and secondary use of patient records. As a result, researchers frequently lack sufficiently large and diverse datasets to develop, test, and compare multimodal Deep Learning architectures in healthcare.

Synthetic data generation has emerged as a possible solution to this limitation. Unlike real-world medical data, synthetic records are not linked to actual patients and can therefore be used for experimentation while reducing privacy concerns. In addition, synthetic data can be generated at scale and under controlled conditions, making it possible to simulate rare but clinically important pathologies that are often underrepresented in real datasets.

This study aims to address the structural weaknesses of purely data-driven generative models, which often output physiologically impossible values, by proposing a highly controlled expert-driven simulation engine. The main contributions of this work are threefold:
\begin{itemize}
    \item Multimodal Synthetic Generation: We engineered a transparent, rule-based physiological pipeline utilizing truncated normal distributions to generate 500,000 synthetic patient records across 24 diagnostic classes, mathematically enforcing clinical coherence.
    \item Real-World Distributional Validation: We validated the clinical representativeness of our synthetic engine by conducting a distributional comparison against a real-world public Emergency Department dataset, ensuring epidemiological realism.
    \item Comprehensive Benchmarking \& XAI Validation: We benchmarked Deep Learning models (Early Fusion MLP and Tabular Transformer) against an established tabular baseline (XGBoost), and used SHAP to verify whether the trained models rely on clinically meaningful features consistent with the generation rules.
\end{itemize}

The remainder of this paper is organized as follows. Section II presents related work, Section III describes the proposed method, Section IV reports the experimental results, and Section V concludes the paper.